\section*{第 \ref{ch:g6:lines} 章 --- 直线、圆与角}

\begin{solution}{exo:g6:lines:1}
自由作图。线段的端点为 $A$ 和 $B$；射线从
$C$ 出发并经过 $A$ 继续延伸；直线在 $B$ 和 $C$ 两侧
均可无限延伸。
\end{solution}

\begin{solution}{exo:g6:lines:2}
“$[AB] = [BA]$”：\emph{正确} --- 两点之间的
部分与点的顺序无关。

“$[AB) = [BA)$”：\emph{错误} --- $[AB)$ 从 $A$ 出发，$[BA)$ 从
$B$ 出发；两者方向相反。

“$(AB) = (BA)$”：\emph{正确} --- 两种写法表示的是同一条
无限延伸的直线。
\end{solution}

\begin{solution}{exo:g6:lines:3}
步骤：沿 $d$ 平移三角尺，直到它的直角边
经过 $P$；画出这条垂线，记为 $p$。然后
用同样的方法画出与 $p$ 垂直且经过 $P$ 的直线，
由\cref{prop:g6:lines:facts}可知，它与 $d$ 平行。
\end{solution}

\begin{solution}{exo:g6:lines:4}
$a$ 与 $b$ 都垂直于同一条直线 $c$，因此
$a \parallel b$（性质 1）。对于 $e$ 与 $c$：二者都
垂直于同一条直线 $a$（已知 $e \perp a$，且 $a \perp c$ 表示
$c \perp a$ 也成立），于是再次应用性质 1，得 $e \parallel c$。
\end{solution}

\begin{solution}{exo:g6:lines:5}
$OM = 3$ cm 恰好（$M$ 在圆上）；$ON < 3$ cm（$N$ 在圆内）；
$OP > 3$ cm（$P$ 在圆外）。
\end{solution}

\begin{solution}{exo:g6:lines:6}
半径 $= 9 \div 2 = 4.5$ cm。\quad 直径 $= 2 \times 2.6 = 5.2$ cm。
\end{solution}

\begin{solution}{exo:g6:lines:7}
$\widehat{SRT} = 72^\circ$：锐角。\quad
$\widehat{ABC} = 148^\circ$：钝角。\quad
$\widehat{MON} = 90^\circ$：直角。（角的顶点总是中间
那个字母。）
\end{solution}

\begin{solution}{exo:g6:lines:8}
度数取决于所画的三角形，但三个角之和
总是（非常接近）$180^\circ$ --- 微小差异来自
测量误差。在\Cref{ch:g7:triangles}中证明了这个和
恰为 $180^\circ$。
\end{solution}

\begin{solution}{exo:g6:lines:9}
自由作图。要在不用量角器的条件下作出 $130^\circ$，可从 $65^\circ$ 出发：
把 $65^\circ$ 的角并排复制两次（$65 + 65 =
130$），例如用透明纸复制，
或用圆规和直尺复制角。
\end{solution}

\begin{solution}{exo:g6:lines:10}
距离为 $3$ cm 的点（即到 $A$ 的距离）构成圆心为 $A$、
半径为 $3$ cm 的圆；距离为 $2$ cm 的点（即到 $B$ 的距离）
构成圆心为 $B$、半径为 $2$ cm 的圆。所求的点即为这两个圆的交点：两圆
相交时（$AB$ 严格介于 $1$ 与
$5$ cm 之间）有两个交点；两圆相切（$AB = 5$ cm 或 $AB = 1$ cm）时有一个交点；两圆
相离或一圆在另一圆内时，没有交点。
\end{solution}

\begin{solution}{exo:g6:lines:11}
$12$ 个整点刻度把整周分成 $12$ 个角，每个角
为 $360 \div 12 = 30^\circ$。3 点时，两针之间跨过 $3$ 个刻度：
$90^\circ$（直角）。5 点时：$5 \times 30 = 150^\circ$。
6 点时：$180^\circ$（平角）。
\end{solution}

\begin{solution}{exo:g6:lines:12}
以 $6$ cm 为半径，分别以 $A$ 和 $B$ 为圆心画两段圆弧，它们
交点为 $C$。三条边长都是 $6$ cm，这个三角形是
等边三角形，每个角都是 $60^\circ$（量得的结果
也印证了这一点）。猜想：等边三角形的三个角都相等，且各为
$60^\circ$。
\end{solution}

\begin{solution}{pb:g6:lines:1}
\textbf{1.} 半周：$360 \div 2 = 180^\circ$。四分之一周：
$360 \div 4 = 90^\circ$。十二分之一周：$360 \div 12 = 30^\circ$。

\textbf{2.} 十二个相等的角恰好填满 $360^\circ$：
相邻刻度之间的角为 $30^\circ$。3 点时，两针
之间跨过 $3$ 个刻度：$3 \times 30 = 90^\circ$ --- 直角，
与\cref{exo:g6:lines:11}中所得结果一致。

\textbf{3.} 1 点时：$1 \times 30 = 30^\circ$。
4 点时：$4 \times 30 = 120^\circ$。7 点时，两针
在一侧跨过 $7$ 个刻度，$7 \times 30 = 210^\circ$，因此
两针之间的角是另一侧的角：
$360 - 210 = 150^\circ$。

\textbf{4.} $360^\circ$ 对应 $60$ 分钟：$360 \div 60 =
6^\circ$ 每分钟。

\textbf{5.} $30^\circ$ 对应 $60$ 分钟：$30 \div 60 = 0.5^\circ$
每分钟 --- 即半度。时针缓慢爬行，分针
则大步前进。

\textbf{6.} 自 3:00 起的 $30$ 分钟内，时针已
走过 $30 \times 0.5 = 15^\circ$，越过了 $3$：它恰好位于
$3$ 与 $4$ 之间的正中间，所成的角为 $90 + 15 = 105^\circ$
（自 $12$ 算起）。分针指向 $6$：
$180^\circ$。两针之间的角为：
$180 - 105 = 75^\circ$。

\textbf{7.} 计算结果推翻了猜测：分针指向
$180^\circ$，但时针已经离开了 $6$ ——它位于
$6 \times 30 + 30 \times 0.5 = 180 + 15 = 195^\circ$。两针之间的角为
$195 - 180 = 15^\circ$：两针很接近，
但并未完全重合。

\textbf{8.} 分针指向 $3$：$90^\circ$。时针：
$9 \times 30 + 15 \times 0.5 = 270 + 7.5 = 277.5^\circ$。
两者之差为 $277.5 - 90 = 187.5^\circ$，因此两针之间的角为
$360 - 187.5 = 172.5^\circ$ ——
几乎是一个平角。

\textbf{9.} 1:00 时，分针（$0^\circ$）\emph{落后}于
时针（$30^\circ$）。1:10 时，分针
（$60^\circ$）\emph{超前}于它
（$30 + 10 \times 0.5 = 35^\circ$）。分针转得更快，
每分钟 $6^\circ$（时针只有 $0.5^\circ$），所以在 1:00 与
1:10 之间，分针追上并超过了时针——
超过的那一刻，两针恰好完全重合。

\textbf{10.} 同样的追赶一再发生：两针
在 12:00 恰好相遇，之后约在 1:05、2:11、3:16、
4:22、5:27、6:33、7:38、8:44、9:49 和 10:55 相遇。人们可能
以为 11 点到 12 点之间还有一次相遇，其实它正好\emph{落在} 12:00，
而这一刻已经开始了下一轮 ——因此在 12 小时内，两针
重合 $11$ 次，而不是 $12$ 次。一整天则重合 $22$ 次。

\textbf{11.} 有公共顶点、既不重叠也无空隙的角，
恰好铺满该顶点周围的一整圈：绕着该点旋转一周，
每个角恰好被扫过一次，而一整圈扫过的角为
$360^\circ$。因此各角的度数之和为 $360^\circ$ ——
与角的个数及大小无关。

\textbf{12.} 一天中每小时对应
$360 \div 24 = 15^\circ$。于是：
\[
\text{睡眠 } 9 \times 15 = 135^\circ, \quad
\text{上学 } 90^\circ, \quad
\text{玩耍 } 45^\circ, \quad
\text{用餐 } 30^\circ, \quad
\text{其他 } 60^\circ,
\]
且 $135 + 90 + 45 + 30 + 60 = 360^\circ$：整个圆饼恰好分完，
既无空隙也无重叠（见第 11 题）。

\textbf{13.} 用圆规画一个圆，并画一条半径（即
起始线）。把量角器的中心放在圆心上，
使其 $0^\circ$ 刻线落在该半径上，并标出
$135^\circ$；画出新的半径，睡眠扇形就完成了。然后
把 $0^\circ$ 刻线放在\emph{刚画出的}那条半径上，
标出上学的 $90^\circ$，依此类推 ——每个扇形都
从前一个扇形结束处开始。画完最后一个扇形后，图形
恰好在起始半径上闭合。

\textbf{14.} $90^\circ$ 占整个圆盘的 $\frac{90}{360} = \frac14$
（见\cref{met:g6:fractions:simplify}），因此它代表一天的四分之一：
$24 \div 4 = 6$ 小时。

\textbf{15.} 时针 $12$ 小时转一圈，一天转 $2$
圈。分针每小时转一圈，共转 $24$ 圈。
秒针每分钟转一圈，一天共转 $24 \times 60 = 1\,440$
圈。
\end{solution}
